\begin{reviewcomment}{1}
	The security implications of the unbound program length should be more strongly addressed.
	According to the manuscript, the preimage of the commitment does not include the length of the program $\ell$.
	Thus, a TPS that deviates from the prescribed workflow may consider a stack-consistent prefix of the committed token array instead of the full desired policy.
	This is not an implementation detail, but a matter of soundness.
	The proposed framework is oriented toward verifiable consent enforcement and therefore the authors should either bind the program length in the commitment/proof relation in the current implementation or give a much stronger justification for postponing this modification.
\end{reviewcomment}

\begin{reviewresponse}{1}
	We agree and have bound the program length in the current implementation.
	The commitment is now $h=\mathsf{Poseidon}(\mathsf{H}_{\mathrm{pol}}(P_{\mathrm{enc}})\parallel\ell\parallel r)$, and the circuit supplies the same private $\ell$ to both the commitment check and the interpreter, so an accepting proof can no longer evaluate a shorter stack-consistent prefix of the committed program.
	The change adds $88$ constraints (about $0.1\%$ of the circuit) and no public input, as Section~6.9 now reports.
	The statement that this binding was left to future work has been removed.

	\RevisionLocation{Section~5.3, first paragraph (\Loc{ds-len}); Section~6.9, paragraph introducing Table~15 (\Loc{ab-bind}); Section~7, first paragraph under ``Policy and context binding'' (\Loc{di-len}).}
\end{reviewresponse}
